% !PRINTED 497
\weberpage{497}
%<<CONT
recognises, and from this it follows that a cyclic permutation belongs to the first or to the second kind, according as the number of figures is odd or even.

Just as every permutation can be composed of transpositions, so too every permutation of the first kind can be composed of three-term cycles. It suffices if this is proved for every pair of transpositions, since every permutation of the first kind can be compounded of such pairs. But we have
\[
(0, 1)(0, 2) = (0, 1, 2)
\]
\[
(0, 3)(1, 2) = (0, 1, 2)(0, 1, 3),
\]
from which the correctness of the assertion is to be seen, since the two transpositions of a pair have either one figure in common or none. Hence:

\gesperrt{6. All permutations of the first kind can be decomposed (in infinitely many ways) into cyclic permutations of three elements.}

From 6. there follows further

\gesperrt{7. A permutation group which contains all three-term cyclic permutations with two fixed figures must contain the whole alternating group.}

For from the cycles $(0, 1, 2), (0, 1, 3), (0, 1, 4) \dots (0, 1, m - 1)$ all three-term cycles, and hence by 6. all permutations of the first kind, can be compounded, as from the compositions
\[
(1, 0, 2) = (0, 1, 2)(0, 1, 2),
\]
\[
(0, 2, 3) = (0, 1, 2)(1, 0, 3)(1, 0, 2),
\]
\[
(1, 3, 4) = (1, 0, 3)(0, 1, 4)(0, 1, 3)
\]
\[
(2, 3, 4) = (1, 0, 2)(1, 3, 4)(0, 1, 2)
\]
may be seen.

\gesperrt{8. If $m$ is greater than 4, then a permutation group which contains all permutations $(0, 1)(2, 3)$ composed of any two transpositions without a common figure is divisible by the alternating group.}

For we have
\[
(0, 1, 2) = (0, 1)(3, 4)(3, 4)(0, 2)
\]
%CONT>>
